
\documentclass[10pt,a4paper]{article}

\usepackage[utf8]{inputenc}
\usepackage[T1]{fontenc}
\usepackage[french]{babel}
\usepackage{amsmath,amssymb}
\usepackage{geometry}
\usepackage{enumitem}

\geometry{
    top=1.5cm,
    bottom=1.5cm,
    left=1.7cm,
    right=1.7cm
}

\setlength{\parindent}{0pt}

\newcommand{\version}[1]{
    \begin{center}
        \textbf{\large INTERROGATION -- LOIS BINOMIALE ET NORMALE}\\
        \textbf{Version #1}
    \end{center}

    \vspace{-0.1cm}

    \textbf{Nom :} \rule{5cm}{0.4pt}
    \hfill
    \textbf{Prénom :} \rule{4cm}{0.4pt}

    \vspace{0.2cm}
}

\begin{document}

% =========================================================
% VERSION A
% =========================================================

\version{A}

\textbf{Exercice 1 -- Loi binomiale \hfill /2}

Dans un laboratoire de microbiologie, on analyse des prélèvements
afin de rechercher la présence d'une bactérie. On estime que la
probabilité qu'un prélèvement soit positif est de $0,08$.
Les résultats obtenus pour les différents prélèvements sont supposés
indépendants.

On analyse successivement $20$ prélèvements et on note $X$ le nombre
de prélèvements positifs.

\begin{enumerate}[label=\arabic*.]
    \item Justifier que $X$ suit une loi binomiale et préciser ses paramètres.
\end{enumerate}


\textbf{Exercice 2 -- Loi normale \hfill /8}

Dans un laboratoire d'analyses biologiques, on mesure la concentration
(en mg/L) d'une molécule dans le sang. Cette concentration est
modélisée par une variable aléatoire $Y$ suivant une loi normale de
moyenne
\[
\mu=100
\]
et d'écart-type
\[
\sigma=12.
\]

\begin{enumerate}[label=\arabic*.]
    \item Rappeler à quoi correspondent les paramètres $\mu$ et $\sigma$.

    \item Calculer
    \[
    P(88\leq Y\leq112).
    \]

    \item Quelle est la probabilité qu'un échantillon présente une
    concentration inférieure à $124$ mg/L ?

    \item Autour de quelles valeurs doit être située la concentration
    pour que $95\,\%$ des échantillons aient une concentration comprise
    entre ces deux valeurs ?

        \textit{Rappels : Si $X \sim \mathcal{N}(m\,;\,\sigma)$, alors, quelle que soit la loi normale considérée~:
$$P(m-\sigma \leqslant X \leqslant m+\sigma) \approx 0{,}68 \qquad P(m-2\sigma \leqslant X \leqslant m+2\sigma) \approx 0{,}95 \qquad P(m-3\sigma \leqslant X \leqslant m+3\sigma) \approx 0{,}997$$} 
\end{enumerate}

\vfill

% =========================================================
% VERSION B
% =========================================================

\version{B}

\textbf{Exercice 1 -- Loi binomiale \hfill /2}

Dans un laboratoire de contrôle microbiologique, on teste des
échantillons d'eau afin de détecter la présence d'un microorganisme.
On estime que la probabilité qu'un échantillon soit positif est de
$0,10$. Les résultats des différents tests sont supposés indépendants.

On réalise $20$ tests et on note $X$ le nombre de tests donnant un
résultat positif.

\begin{enumerate}[label=\arabic*.]
    \item Justifier que $X$ suit une loi binomiale et préciser ses paramètres.
\end{enumerate}



\textbf{Exercice 2 -- Loi normale \hfill /8}

Dans un laboratoire de biochimie, on mesure la concentration (en
mmol/L) d'une substance dans des échantillons sanguins. Cette
concentration est modélisée par une variable aléatoire $Y$ suivant
une loi normale de moyenne
\[
\mu=5,0
\]
et d'écart-type
\[
\sigma=0,6.
\]

\begin{enumerate}[label=\arabic*.]
    \item Rappeler à quoi correspondent les paramètres $\mu$ et $\sigma$.

    \item Calculer
    \[
    P(4,4\leq Y\leq5,6).
    \]

    \item Quelle est la probabilité qu'un échantillon présente une
    concentration inférieure à $6,2$ mmol/L ?

    \item Autour de quelles valeurs doit être située la concentration
    pour que $99,7\,\%$ des échantillons aient une concentration comprise
    entre ces deux valeurs ?

        \textit{Rappels : Si $X \sim \mathcal{N}(m\,;\,\sigma)$, alors, quelle que soit la loi normale considérée~:
$$P(m-\sigma \leqslant X \leqslant m+\sigma) \approx 0{,}68 \qquad P(m-2\sigma \leqslant X \leqslant m+2\sigma) \approx 0{,}95 \qquad P(m-3\sigma \leqslant X \leqslant m+3\sigma) \approx 0{,}997$$} 
\end{enumerate}

\newpage

% =========================================================
% VERSION C
% =========================================================

\version{C}

\textbf{Exercice 1 -- Loi binomiale \hfill /2}

Dans un laboratoire de biologie médicale, on recherche une substance
particulière dans des échantillons de sang. On estime que la probabilité
qu'un échantillon contienne cette substance est de $0,15$.
Les analyses sont supposées indépendantes.

On analyse $20$ échantillons et on note $X$ le nombre d'échantillons
dans lesquels la substance est détectée.

\begin{enumerate}[label=\arabic*.]
    \item Justifier que $X$ suit une loi binomiale et préciser ses paramètres.
\end{enumerate}



\textbf{Exercice 2 -- Loi normale \hfill /8}

Dans un laboratoire pharmaceutique, on mesure la masse (en mg) de
principe actif contenue dans des comprimés. Cette masse est modélisée
par une variable aléatoire $Y$ suivant une loi normale de moyenne
\[
\mu=250
\]
et d'écart-type
\[
\sigma=10.
\]

\begin{enumerate}[label=\arabic*.]
    \item Rappeler à quoi correspondent les paramètres $\mu$ et $\sigma$.

    \item Calculer
    \[
    P(230\leq Y\leq270).
    \]

    \item Quelle est la probabilité qu'un comprimé contienne moins de
    $280$ mg de principe actif ?

    \item Autour de quelles valeurs doit être située la masse de
    principe actif pour que $95\,\%$ des comprimés aient une masse
    comprise entre ces deux valeurs ?

        \textit{Rappels : Si $X \sim \mathcal{N}(m\,;\,\sigma)$, alors, quelle que soit la loi normale considérée~:
$$P(m-\sigma \leqslant X \leqslant m+\sigma) \approx 0{,}68 \qquad P(m-2\sigma \leqslant X \leqslant m+2\sigma) \approx 0{,}95 \qquad P(m-3\sigma \leqslant X \leqslant m+3\sigma) \approx 0{,}997$$} 
\end{enumerate}

\vfill

% =========================================================
% VERSION D
% =========================================================

\version{D}

\textbf{Exercice 1 -- Loi binomiale \hfill /2}

Dans un laboratoire de contrôle qualité, on recherche la présence
d'une contamination bactérienne dans des prélèvements issus d'une
production. On estime que la probabilité qu'un prélèvement soit
contaminé est de $0,05$. Les résultats des analyses sont supposés
indépendants.

On réalise $20$ analyses et on note $X$ le nombre de prélèvements
contaminés.

\begin{enumerate}[label=\arabic*.]
    \item Justifier que $X$ suit une loi binomiale et préciser ses paramètres.
\end{enumerate}

\vfill

\textbf{Exercice 2 -- Loi normale \hfill /8}

Dans un laboratoire d'analyses médicales, on mesure la concentration
(en g/L) d'une protéine dans des échantillons de sang. Cette
concentration est modélisée par une variable aléatoire $Y$ suivant
une loi normale de moyenne
\[
\mu=40
\]
et d'écart-type
\[
\sigma=4.
\]

\begin{enumerate}[label=\arabic*.]
    \item Rappeler à quoi correspondent les paramètres $\mu$ et $\sigma$.

    \item Calculer
    \[
    P(36\leq Y\leq44).
    \]

    \item Quelle est la probabilité qu'un échantillon présente une
    concentration inférieure à $48$ g/L ?

    \item Autour de quelles valeurs doit être située la concentration
    pour que $99,7\,\%$ des échantillons aient une concentration comprise
    entre ces deux valeurs ?

    \textit{Rappels : Si $X \sim \mathcal{N}(m\,;\,\sigma)$, alors, quelle que soit la loi normale considérée~:
$$P(m-\sigma \leqslant X \leqslant m+\sigma) \approx 0{,}68 \qquad P(m-2\sigma \leqslant X \leqslant m+2\sigma) \approx 0{,}95 \qquad P(m-3\sigma \leqslant X \leqslant m+3\sigma) \approx 0{,}997$$} 
\end{enumerate}

\end{document}

